% ============================================================
% Autor: Eloi Egea Rada
% Tema: Development, Analysis, and Results of the Bachelor's Thesis Project
% ============================================================

\chapter{Development}

This chapter describes the implementation of the penetration testing lab environment
built on EVE-NG. It covers the platform setup, the management of virtual machine
images, the lab management interface, the infrastructure problems encountered during
development, and an overview of the lab topologies designed for the course. Each
individual lab is then described in its own section, with a summary of the network
topology, techniques, and problems encountered. Detailed instructor references ---
including deployment credentials, startup procedures, and flag locations --- are
compiled in Appendices~
ef{app:lab1}--
ef{app:lab3}. The EVE-NG installation
procedure is documented in the \href{https://TheEgea.github.io/TFG/annexos/app-a/}{project online documentation}.


% ------------------------------------------------------------
\section{Platform Setup}
% ------------------------------------------------------------

The development process began with the installation and validation of EVE-NG~\cite{eve_ng} as the
network emulation platform. Following the official EVE-NG documentation, an initial
environment was deployed in order to verify basic functionality: creating simple
network topologies, deploying virtual machines with different operating systems, and
testing console and remote desktop access before committing to the full lab design.

Several attempts were made to run EVE-NG as a virtual machine on a VMware-based
infrastructure available in the initial university environment. Despite VMware being
an officially supported hypervisor, recurrent stability and integration issues
prevented a reliable setup. After evaluating the available alternatives and in
agreement with the project supervisor, the platform was migrated to a dedicated
Proxmox VE~\cite{proxmox} server. This choice offers better compatibility with the nested
virtualisation required by EVE-NG and can be replicated on any institutional server
running the same hypervisor.

On top of Proxmox, an EVE-NG virtual machine was deployed following the official
installation guidelines, with the networking model adapted to the Proxmox bridge
configuration. Console access via Telnet and graphical access via VNC were validated
successfully for all node types used in the labs. Native Wireshark integration was
not enabled in this iteration due to conflicts with pre-existing software on the host
machine, and was therefore left out of scope for the current version.


% ------------------------------------------------------------
\section{Image Management}
% ------------------------------------------------------------

\subsection{Image Selection}

Once the platform was stable, a survey of available EVE-NG-compatible images was
conducted, prioritising free and open-source operating systems suitable for an
academic context. Starting from a broad list of candidate images, the final selection
was reduced to a minimal, reusable core for the penetration testing scenarios:

\begin{itemize}
    \item \textbf{VyOS}~\cite{vyos}: used as a router and Layer~3 switch for routing and
          inter-VLAN scenarios.
    \item \textbf{pfSense CE}~\cite{pfsense}: used as the perimeter firewall, providing NAT,
          DNS resolution, and port forwarding.
    \item \textbf{GNU/Linux distributions} (Ubuntu Server~24.04~\cite{ubuntu}, Ubuntu
          Desktop~24.04, Parrot OS~\cite{parrot_os}): used as target servers, internal workstations,
          and attacker machines respectively.
\end{itemize}

This reduced set balances realism and manageability, keeps all software licences
aligned with the academic context, and limits the overhead of maintaining multiple
image versions across labs.

Parrot OS Security Edition was selected over the more widely known Kali
Linux~\cite{kali_linux} after evaluating both options during image integration.
The Kali image proved substantially heavier, required additional QEMU guest
configuration to work reliably inside EVE-NG, and produced higher per-instance
resource overhead. Parrot OS provides an equivalent toolset --- Nmap, Metasploit,
Burp Suite, Hydra, Gobuster, and all tools required by the labs --- with a
lighter footprint and cleaner EVE-NG integration. An additional consideration
was accessibility: Parrot OS includes OpenDyslexic and other accessibility fonts
by default, consistent with the inclusive design principles adopted for the lab
documentation. This decision is reflected in the requirements specification
(Section~TR2.1).


\subsection{Image Organisation in EVE-NG}

To ensure consistency and reproducibility, all images were organised under the
standard EVE-NG path \path{/opt/unetlab/addons/qemu/}. Each system was placed in
a dedicated directory with a normalised name following the pattern
\texttt{<type>-<version>} (e.g., \path{pfsense-2.7.2},
\path{vyos-1.x-rolling-amd64}, \path{linux-ubuntu-server-24.04}).

For GNU/Linux systems, the \texttt{linux-} prefix was applied consistently
(e.g., \path{linux-ubuntu-desktop-24.04}, \path{linux-parrot-rolling}), since
EVE-NG only makes the \textit{Linux} node category available in the interface when
at least one directory with this prefix exists. Each directory contains a bootable
image named \path{cdrom.iso} and a \path{virtioa.qcow2} disk image sized
according to the role of the node.

\subsection{ISO Upload Automation}

The manual workflow for preparing each image --- creating the directory, uploading
the ISO, renaming it, creating the disk, and fixing permissions --- proved repetitive
and error-prone during initial experiments. To address this, a Python script named
\texttt{Iso\_uploader.py} was developed as part of the project tooling.

The script presents a graphical file picker, uploads the selected ISO files to the
EVE-NG server via SCP, and then automates the remaining preparation steps: directory
creation under \path{/opt/unetlab/addons/qemu/}, renaming the ISO to
\path{cdrom.iso}, creation of the \path{virtioa.qcow2} disk at a size inferred
from the image type, and execution of \texttt{unl\_wrapper -a fixpermissions}. This
tool reduces the time required to onboard a new image and enforces a uniform
directory structure across deployments.


% ------------------------------------------------------------
\section{EVE-NG Lab Management Interface}
% ------------------------------------------------------------

EVE-NG provides a web-based graphical interface through which administrators and
students can manage lab files independently of each other. The \textit{File Manager},
accessible after login, is the main entry point for lab operations and offers several
features that are relevant to reproducibility and multi-user deployments:

\begin{itemize}
    \item \textbf{Lab import and export}: Topology files (\texttt{.unl}) can be
          imported by dragging them into the interface. Exported archives preserve
          the topology, node configurations, and snapshots, making it straightforward
          to share or back up a complete lab.
    \item \textbf{Folder structure}: Labs can be organised hierarchically
          (e.g., \texttt{/GEISI-Labs/Pentesting/}), which keeps different courses
          and iterations separate without requiring multiple EVE-NG installations.
    \item \textbf{Multi-lab support}: Several topologies can run simultaneously.
          The administrator dashboard provides a real-time view of CPU and memory
          consumption per lab, which is useful for capacity planning on shared
          infrastructure.
\end{itemize}

\begin{figure}[h]
    \centering
    \includegraphics[width=0.9\textwidth]{eve-file-manager.png}
    \caption{EVE-NG File Manager showing the lab directory structure, import
             controls, and topology thumbnails.}
    \label{fig:eve-file-manager}
\end{figure}

The export format also simplifies distribution: a \texttt{.unl} file committed to
the project repository can be imported into any compatible EVE-NG instance,
reproducing the full topology without manual reconfiguration.


% ------------------------------------------------------------
\section{Infrastructure Problems}
% ------------------------------------------------------------

Several problems were encountered at the infrastructure level during the setup of
the EVE-NG environment. These are distinct from the lab-specific issues documented
within each lab section and relate to the underlying platform and tooling.

\subsection{EVE-NG on VMware}

The initial deployment target was a VMware-based server available through the
university. Despite EVE-NG being officially listed as compatible with VMware~\cite{vmware_esxi}
hypervisors, the environment produced recurrent instability: nodes failed to start
consistently, VNC sessions dropped without warning, and network connectivity between
virtual bridges was unreliable. No stable configuration was found after several
iterations. The root cause was not fully identified, but the combination of
VMware's nested virtualisation restrictions and the specific hardware and software
versions in the university environment appeared to be contributing factors.
The decision to migrate to Proxmox resolved all of these issues.

\subsection{Node Configuration Persistence}

A recurring challenge in EVE-NG lab design is ensuring that node configurations
survive restarts. By default, runtime network settings applied inside guest virtual
machines are lost when the lab is stopped or the EVE-NG host is rebooted. This
affected multiple nodes across the lab and required a dedicated solution for each
one. The specific mechanisms used are described in the relevant lab sections; the
general principle adopted throughout the project was to rely on the native
persistence mechanisms of each operating system --- such as \texttt{systemd-networkd}
unit files, NetworkManager connection profiles, and \texttt{udev} rules --- rather
than applying settings interactively at the console.

\subsection{LVM Volume Group Naming Conflict}

During a session involving disk operations on guest virtual machine images mounted
from the EVE-NG host, a naming conflict arose between the host's LVM volume group
and a guest disk's volume group, both named \texttt{ubuntu-vg}. An attempt to
resolve the conflict by renaming the host volume group caused the EVE-NG host to
fail on the next restart, dropping into an \texttt{initramfs} prompt because the
boot configuration still referenced the original group name.

Recovery was performed from the \texttt{initramfs} shell by renaming the volume
group back to its original name and reactivating it before resuming the normal boot
sequence. As a result of this incident, the rule established for the project is to
never rename a volume group belonging to the host operating system. When a naming
conflict arises with a guest disk, only the guest volume group is renamed, using a
descriptive identifier that includes the lab and node reference.

\subsection{Memory Capacity Constraint}

During concurrent development of the first two laboratories, peak memory
consumption on the EVE-NG virtual machine regularly exceeded the initial
32\,GB RAM allocation of the host server. Running pfSense, VyOS, Ubuntu Server,
Ubuntu Desktop, and Parrot OS simultaneously --- alongside other services
co-hosted on the same Proxmox infrastructure (monitoring stack, reverse proxy,
storage services) --- produced memory pressure that caused nodes to be killed
by the hypervisor's out-of-memory handler, interrupting development sessions.

An additional 32\,GB RAM module was acquired and installed, bringing the host
to 64\,GB total. This resolved the stability issues and allowed all lab nodes
to run concurrently without further intervention. The upgrade was not included
in the original project budget estimate; its approximate cost
(\textasciitilde 80\,\euro) is noted in Chapter~6 as a post-planning expenditure.


% ------------------------------------------------------------
\section{Lab Topologies}
% ------------------------------------------------------------

Four penetration testing lab topologies were designed for the course. Each topology
is implemented as an EVE-NG \texttt{.unl} file and can be imported into any
compatible EVE-NG instance. Table~\ref{tab:lab-overview} summarises the four labs,
their primary learning objective, and their current implementation status.

\begin{table}[h]
\centering
\small
\setlength{\tabcolsep}{5pt}
\begin{tabularx}{\textwidth}{@{}lXl@{}}
\toprule
\textbf{Lab} & \textbf{Primary Objective} & \textbf{Status} \\
\midrule
Lab~1: Reconnaissance \& Enumeration & Port scanning, web content discovery, SSH & Implemented \\
Lab~2: Web Application Vulnerabilities & OWASP Top~10 exploitation & Implemented \\
Lab~3: Network Analysis & Traffic capture, lateral movement, incident response & Implemented \\
Lab~4: Privilege Escalation & Local privilege escalation techniques & Pending \\
\bottomrule
\end{tabularx}
\caption{Overview of the four EVE-NG lab topologies designed for the course.}
\label{tab:lab-overview}
\end{table}

The remainder of this chapter describes each implemented lab in detail, covering
the network topology, design decisions, student workflow, and problems encountered
during development.


% ------------------------------------------------------------
% Lab descriptions --- one file per lab
% ------------------------------------------------------------
% Lab 1 --- compact version (full docs: docs/web/docs/labs/lab1/)
\section{Lab 1: Reconnaissance and Enumeration}

Lab~1 introduces the reconnaissance phase of a penetration test following the PTES
framework~\cite{ptes}. Students interact with a fictional target organisation,
\textit{PEBCAK Corp}, through a layered topology that includes a perimeter firewall
(pfSense~CE~\cite{pfsense}), an internal router (VyOS~\cite{vyos}), and a target
Ubuntu~Server~24.04~\cite{ubuntu}. Access is exclusively via DNS hostname
(\texttt{lab1}) resolved through pfSense's Unbound resolver~\cite{unbound_dns}, requiring students
to use proper service enumeration before any exploitation attempt. A teacher reference --- credentials, network table, and startup checklist --- is compiled
in Appendix~B; the complete topology, configuration files, and step-by-step guide are at
\path{docs/web/docs/labs/lab1/}.


\begin{figure}[H]
\centering
\includegraphics[width=0.95\textwidth]{lab1-topology.png}
\caption{Lab~1 network topology}
\label{fig:lab1-topology}
\end{figure}

\newpage
\subsubsection*{Network Addressing}

\begin{table}[H]
\centering
\resizebox{\textwidth}{!}{\begin{tabular}{@{}lllll@{}}
\toprule
\textbf{Node} & \textbf{Role} & \textbf{Interface} & \textbf{IP} & \textbf{Network} \\
\midrule
Parrot-OS   & Attacker        & eth0 & 192.168.0.x (DHCP) & 192.168.0.0/24 (Homelab)  \\
pfSense     & Firewall/NAT    & WAN  & 192.168.0.29       & 192.168.0.0/24            \\
pfSense     & Firewall/NAT    & LAN  & 172.16.0.1         & 172.16.0.0/30             \\
VyOS        & Router/L3       & eth0 & 172.16.0.2         & 172.16.0.0/30             \\
VyOS        & Router/L3       & eth6 & 192.168.10.5       & 192.168.10.0/24 (Users)   \\
VyOS        & Router/L3       & eth7 & 192.168.20.1       & 192.168.20.0/24 (Servers) \\
Server      & Target (web+SSH)& eth0 & 192.168.20.50      & 192.168.20.0/24           \\
PC1         & Workstation     & eth0 & 192.168.10.50      & 192.168.10.0/24           \\
\bottomrule
\end{tabular}}
\caption{Lab~1 IP addressing scheme}
\label{tab:lab1-ip}
\end{table}

\subsubsection*{Techniques and Framework Alignment}

\begin{table}[H]
\centering
\small
\begin{tabularx}{\textwidth}{lXl}
\toprule
\textbf{Technique} & \textbf{Description} & \textbf{Reference} \\
\midrule
DNS resolution        & Hostname-to-IP via Unbound resolver (RFC~1035~\cite{rfc1035}) & PTES Recon  \\
Port scanning         & TCP SYN scan with Nmap (RFC~793~\cite{rfc793})                & PTES Recon  \\
Service enumeration   & Banner grabbing, version detection                            & PTES Recon  \\
Web content discovery & Directory brute-force with Gobuster                           & OWASP WSTG  \\
Credential retrieval  & Exposed credentials in web content                            & OWASP A05:2021\\
SSH access            & Password-based login (RFC~4252~\cite{rfc4252})                & PTES Exploit\\
\bottomrule
\end{tabularx}
\caption{Lab~1 techniques and framework mapping}
\label{tab:lab1-techniques}
\end{table}

\subsection{Problems Encountered and Solutions}

\subsubsection{VNC Keyboard Input on QEMU Nodes}

EVE-NG provides VNC console access to QEMU nodes, but certain key combinations ---
most notably \texttt{Ctrl+C} and special characters such as \texttt{|}, \texttt{@},
and \texttt{\{} --- are not reliably forwarded by the browser-based VNC client or by
the \texttt{vncdo} command-line tool. The root cause was that vncdo sent key symbols
rather than raw key codes, and the VNC server's key mapping did not cover all
combinations used in a penetration testing context.

The solution was to inject keystrokes directly via the QEMU monitor socket,
using Python's raw socket interface to send \texttt{sendkey} commands. This bypasses
the VNC protocol entirely and ensures that arbitrary key sequences, including those
involving modifier keys, are delivered correctly to the guest.

\subsubsection{pfSense DNS Resolver WAN Interface}

After initial configuration, DNS queries from the attacker machine were not being
resolved by pfSense's Unbound resolver. The problem was that Unbound, by default,
only listens on the LAN interface, so queries arriving on the WAN-side interface
were silently discarded. The fix was to explicitly add the WAN interface to
Unbound's listening configuration under \textit{Services $\to$ DNS Resolver $\to$
Network Interfaces}.

\subsubsection{Unbound Startup Configuration Path}

On subsequent restarts, the Unbound configuration applied via the
GUI was not being preserved, causing DNS resolution to fail after every reboot.
Investigation revealed that pfSense's startup scripts write the Unbound configuration
to \path{/var/unbound/unbound.conf} at boot time, overwriting any manual edits.
The correct approach is to configure all resolver settings exclusively through the
GUI or through pfSense's \texttt{/etc/inc/} PHP hooks, never by editing the
generated configuration file directly.

\subsubsection{VyOS Source NAT Requirement}

Traffic routed through VyOS from the attacker to the target was being dropped
because the target's default gateway was pfSense, not VyOS. Return packets were
therefore routed back via pfSense rather than through VyOS, breaking the connection.
The fix was to configure Source NAT (masquerade) on VyOS for traffic entering
the internal segment, ensuring that return traffic would follow the same path back
through the VyOS router.

\subsubsection*{Results}

\begin{table}[H]
\centering
\resizebox{\textwidth}{!}{\begin{tabular}{@{}lll@{}}
\toprule
\textbf{Flag} & \textbf{Location} & \textbf{Technique} \\
\midrule
\texttt{FLAG\{p3bc4k\_s3rv3r\_0wn3d\}} & \texttt{\textasciitilde/flag.txt} on Server via SSH & Credentials from hidden web page \\
\texttt{FLAG\{d3f4ult\_cr3ds\_k1ll\_s3cur1ty\}} & \texttt{/root/flag.txt} on pfSense via SSH pivot & Default credentials \\
\texttt{FLAG\{sk1n3t\_b3c4m3\_s3nt13nt\}} & \texttt{\textasciitilde/flag.txt} on PC1 via lateral move & Credentials found in pfSense flag \\
\bottomrule
\end{tabular}}
\caption{Lab~1 flags and acquisition techniques}
\label{tab:lab1-flags}
\end{table}

% Lab 2 --- compact version (full docs: docs/web/docs/labs/lab2/)
\section{Lab 2: Web Application Vulnerabilities --- SYNAPSE Intelligence Portal}

Lab~2 presents students with a custom two-tier web application,
\textit{SYNAPSE Intelligence Portal}, deployed on a segmented DMZ architecture.
The scenario is aligned with the OWASP Top~10~\cite{owasp2021} and chains
multiple vulnerabilities across two servers: Server-A hosts the portal front-end
with authentication and reporting features; Server-B hosts the DataVault internal
file management service accessible only after pivoting from Server-A. Students must
complete a progressive exploitation chain --- stored XSS, cookie theft, SQL injection,
and YAML deserialization --- to collect all three flags. A teacher reference --- credentials, startup procedure, and flag guide --- is compiled in Appendix~C;
complete source code, topology, and deployment instructions are at \path{docs/web/docs/labs/lab2/}.


\begin{figure}[H]
\centering
\includegraphics[width=0.95\textwidth]{lab2-topology.png}
\caption{Lab~2 network topology}
\label{fig:lab2-topology}
\end{figure}

\subsubsection*{Network Addressing}

\begin{table}[H]
\centering
\resizebox{\textwidth}{!}{\begin{tabular}{@{}lllll@{}}
\toprule
\textbf{Node} & \textbf{Role} & \textbf{Interface} & \textbf{IP} & \textbf{Network} \\
\midrule
Parrot-OS    & Attacker          & eth0  & 10.0.40.10       & 10.0.40.0/24 (Net-Attackers) \\
pfSense-LAB2 & Firewall/NAT      & WAN   & 192.168.0.x (DHCP) & 192.168.0.0/24 (Homelab)   \\
pfSense-LAB2 & Firewall/NAT      & LAN   & 172.16.1.1       & 172.16.1.0/30              \\
VyOS-LAB2    & DMZ Router        & eth0  & 172.16.1.2       & 172.16.1.0/30              \\
VyOS-LAB2    & DMZ Router        & eth1  & 192.168.30.1     & 192.168.30.0/24 (Net-DMZ)  \\
VyOS-LAB2    & DMZ Router        & eth2  & 10.0.40.1        & 10.0.40.0/24 (Net-Attackers)\\
Server-A     & SYNAPSE Portal    & eth0  & 192.168.30.10    & 192.168.30.0/24            \\
Server-B     & DataVault         & eth0  & 192.168.30.20    & 192.168.30.0/24            \\
\bottomrule
\end{tabular}}
\caption{Lab~2 IP addressing scheme}
\label{tab:lab2-ip}
\end{table}

\subsubsection*{OWASP Coverage and Vulnerability Chain}

\begin{table}[H]
\centering
\small
\begin{tabularx}{\textwidth}{llXl}
\toprule
\textbf{Step} & \textbf{OWASP} & \textbf{Vulnerability} & \textbf{Flag} \\
\midrule
1 & A03:2021 & Stored XSS in report field $\rightarrow$ session cookie theft & -- \\
2 & A07:2021 & Cookie injection $\rightarrow$ admin panel access & -- \\
3 & A03:2021 & SQL injection in search $\rightarrow$ credential dump & FLAG\{synapse\_sqli\_creds\_dumped\} \\
4 & A01:2021 & Lateral movement to Server-B via stolen creds & FLAG\{synapse\_pivot\_server\_b\} \\
5 & A08:2021 & YAML deserialization in DataVault upload $\rightarrow$ RCE & FLAG\{synapse\_rce\_achieved\} \\
\bottomrule
\end{tabularx}
\caption{Lab~2 OWASP Top~10 vulnerability chain}
\label{tab:lab2-vulns}
\end{table}

\subsubsection*{Results}

All three flags were successfully captured, confirming the complete exploitation
chain from stored XSS to remote code execution.

\begin{table}[H]
\centering
\resizebox{\textwidth}{!}{\begin{tabular}{@{}lll@{}}
\toprule
\textbf{Flag} & \textbf{Location} & \textbf{Technique} \\
\midrule
\texttt{FLAG\{synapse\_sqli\_creds\_dumped\}} & \texttt{/search?q=} UNION dump on Server-A & UNION SQL Injection \\
\texttt{FLAG\{synapse\_pivot\_server\_b\}}    & \texttt{/home/monitor/flag.txt} on Server-B & SSH pivot with SQLi-extracted credentials \\
\texttt{FLAG\{synapse\_rce\_achieved\}}       & \texttt{/preview} reverse shell on Server-B  & YAML deserialization RCE \\
\bottomrule
\end{tabular}}
\caption{Lab~2 flags and acquisition techniques}
\label{tab:lab2-flags}
\end{table}

\newpage

\subsection{Problems Encountered and Solutions}

\subsubsection{pfSense Block Private Networks and Reply-To}

Two pfSense default behaviours blocked traffic in the lab topology. First, the
\textit{Block private networks} option on the WAN interface was enabled by default,
causing pfSense to silently discard packets arriving from the simulated
attacker network (RFC~1918 space~\cite{rfc1918}). Second, the \textit{Reply-to}
feature added return routes in \texttt{pf} rules that conflicted with VyOS's
routing table, causing asymmetric routing. Both were disabled: the private network
block via the WAN interface settings, and reply-to via a per-rule option in the
firewall rule editor.

\subsubsection{VNC Keyboard Input on pfSense}

The same VNC key injection issue documented in Lab~1 reappeared on pfSense
console sessions, particularly when entering characters required by the pfSense
shell (\texttt{|}, \texttt{\&}, \texttt{\$}). The QEMU monitor raw socket approach
was applied identically.

\subsubsection{Server-A Volume Group Conflict}

When mounting the Server-A QCOW2 image from the EVE-NG host to pre-stage
application files, the host's own LVM volume group (\texttt{ubuntu-vg}) conflicted
with the guest disk's identically named group. Activating both simultaneously
caused the host's logical volumes to be shadowed. The solution was to rename the
guest volume group to \texttt{synapse-vg} using \texttt{vgrename} before activating
it on the host, then rename it back after dismounting.

\subsubsection{Server-B Installer Keyboard Layout}

During the Server-B installation, the Ubuntu installer defaulted to a Spanish
keyboard layout due to the locale of the EVE-NG host, causing mismatched characters
when entering the initial credentials. This was corrected by explicitly selecting the
\texttt{en\_US} keyboard layout at the installer prompt and verified with a
\texttt{/dev/urandom}-seeded password test before proceeding.

\subsubsection{Static Route Gateway Reference}

VyOS rejected a static route configuration that referenced a gateway IP not directly
reachable on any configured interface. The intended next-hop was on a subnet not yet
assigned to any VyOS interface at the time the route was applied. This was resolved
by defining the interface address before the static route, and by using
\texttt{commit} and \texttt{save} in the correct order within the VyOS configuration
session.


% Lab 3 --- compact version (full docs: docs/web/docs/labs/lab3/)
\section{Lab 3: Incident Response and Log Forensics --- HELIX Systems}

Lab~3 inverts the student's role: instead of attacking, they act as a security
analyst responding to an active incident at \textit{HELIX Systems}. The lab
environment comes pre-staged with evidence of a simulated intrusion --- brute-force
SSH attempts, a successful login, privilege escalation via \texttt{sudo}
misconfiguration, backdoor account creation, and lateral movement to a database
server. Students must reconstruct the attack timeline from system logs
(\texttt{auth.log}, \texttt{bash\_history}) and identify the attacker's actions
and impact, following a methodology aligned with NIST SP~800-61~\cite{nist_sp80061}.
A quick teacher reference --- account table, attack timeline, and startup checklist --- is compiled in Appendix~D;
the full evidence walkthrough, forensic guide, and deployment instructions are at \path{docs/web/docs/labs/lab3/}.


\begin{figure}[H]
\centering
\includegraphics[width=0.95\textwidth]{lab3-topology.png}
\caption{Lab~3 network topology}
\label{fig:lab3-topology}
\end{figure}

\subsubsection*{Network Addressing}

\begin{table}[H]
\centering
\resizebox{\textwidth}{!}{\begin{tabular}{@{}lllll@{}}
\toprule
\textbf{Node} & \textbf{Role} & \textbf{Interface} & \textbf{IP} & \textbf{Network} \\
\midrule
Parrot-OS    & Analyst station   & eth0  & 192.168.0.x (DHCP) & 192.168.0.0/24 (Homelab)       \\
pfSense-LAB3 & Firewall/DNAT     & WAN   & 192.168.0.x (DHCP) & 192.168.0.0/24                 \\
pfSense-LAB3 & Firewall          & LAN   & 172.16.2.1         & 172.16.2.0/30                  \\
VyOS-LAB3    & Internal router   & eth0  & 172.16.2.2         & 172.16.2.0/30                  \\
VyOS-LAB3    & Internal router   & eth1  & 192.168.50.1       & 192.168.50.0/24 (Net-Servers)  \\
VyOS-LAB3    & Internal router   & eth2  & 192.168.60.1       & 192.168.60.0/24 (Net-Internal) \\
Server-Web   & Compromised web   & eth0  & 192.168.50.10      & 192.168.50.0/24                \\
Server-DB    & Internal DB       & eth0  & 192.168.60.10      & 192.168.60.0/24                \\
\bottomrule
\end{tabular}}
\caption{Lab~3 IP addressing scheme}
\label{tab:lab3-ip}
\end{table}

\subsubsection*{MITRE ATT\&CK Coverage}

\begin{table}[H]
\centering
\small
\begin{tabularx}{\textwidth}{llX}
\toprule
\textbf{MITRE ID} & \textbf{Technique} & \textbf{Evidence in lab} \\
\midrule
T1110.001 & Brute Force: Password Guessing & \texttt{auth.log} failed attempts \\
T1078     & Valid Accounts                 & \texttt{devops} account compromise \\
T1548.003 & Abuse Elevation Control Mechanism: Sudo and Sudo Caching          & NOPASSWD rule for \texttt{find} in \texttt{sudoers.d/devops} \\
T1136.001 & Create Account: Local          & Backdoor \texttt{maint} account \\
T1021.004 & Remote Services: SSH           & Lateral movement to Server-DB \\
T1005     & Data from Local System         & \texttt{/opt/helix/data/} exfil marker \\
\bottomrule
\end{tabularx}
\caption{Lab~3 MITRE ATT\&CK technique coverage}
\label{tab:lab3-mitre}
\end{table}

\newpage
\subsubsection*{Forensic Findings}

\begin{table}[H]
\centering
\small
\resizebox{\textwidth}{!}{\begin{tabular}{@{}lll@{}}
\toprule
\textbf{Finding} & \textbf{Artefact / Location} & \textbf{Host} \\
\midrule
Brute-force source (\texttt{185.220.101.47}) & \texttt{/var/log/auth.log} --- 25 failed SSH for \texttt{devops} & Server-Web \\
Initial compromise (03:17)   & \texttt{auth.log} successful login as \texttt{devops}                                        & Server-Web \\
Privilege escalation (03:18) & \texttt{sudo /usr/bin/find} NOPASSWD; \texttt{auth.log} + \texttt{/etc/sudoers.d/devops}    & Server-Web \\
Backdoor persistence (03:19) & \texttt{useradd maint}; \texttt{/etc/passwd} + \texttt{/etc/sudoers.d/maint}                 & Server-Web \\
Attacker commands            & \texttt{/home/devops/.bash\_history}                                                         & Server-Web \\
Lateral movement (03:21)     & \texttt{auth.log} \texttt{devops} SSH from \texttt{192.168.50.10}                            & Server-DB  \\
Data exfiltration (03:22)    & \texttt{/opt/helix/data/.exfil\_marker} + \texttt{clients.db}                               & Server-DB  \\
\bottomrule
\end{tabular}}
\caption{Lab~3 forensic findings and evidence locations (all timestamps UTC as recorded in \texttt{/var/log/auth.log})}
\label{tab:lab3-findings}
\end{table}

\subsection{Problems Encountered and Solutions}

\subsubsection{LVM Volume Group Naming Conflict}

When pre-staging Server-Web's evidence files from the EVE-NG host, the same
volume group conflict from Lab~2
recurred: both the host and the guest disk used \texttt{ubuntu-vg}. The resolution
procedure was identical --- renaming the guest VG to \texttt{helix-web-vg} before
activation --- and the naming rule was formally added to the project's infrastructure
runbook to prevent recurrence in future labs.

\subsubsection{Read-only NBD Device}

On one occasion, the QCOW2 image mounted via \texttt{qemu-nbd} was presented as
read-only despite the absence of the \texttt{--read-only} flag. Investigation revealed that a
previous mount session had not been cleanly disconnected (\texttt{qemu-nbd~-{}-disconnect}), leaving the kernel's NBD client in an inconsistent state.
The fix was to explicitly disconnect any stale NBD device before re-mounting, and
to add a pre-mount check to the project's image manipulation scripts.

\subsubsection{Analyst Group Membership for Log Access}

The student account (\texttt{analyst}) was created with restricted permissions so
that it could not read system logs directly without elevated group membership. However, \texttt{/var/log/auth.log}
requires membership of the \texttt{adm} group on Ubuntu. The pre-staging script was
updated to add \texttt{analyst} to \texttt{adm} at image preparation time, preserving
the educational constraint (no \texttt{sudo} to \texttt{root}) while allowing access
to the evidence files that are the focus of the lab.

\subsubsection{pfSense Console Access via QEMU Monitor}

The pfSense serial console was inaccessible through the standard EVE-NG Telnet
interface due to a configuration mismatch between the QEMU serial port definition
and pfSense's expected console device (\texttt{ttyu0} vs \texttt{ttyS0}). Console
access was recovered by connecting directly to the QEMU monitor socket and issuing
a \texttt{sendkey} sequence to navigate the pfSense boot menu, where the console
port assignment could be corrected persistently.
